\documentclass[11pt,a4paper]{report}
\usepackage[T1]{fontenc}    % Indispensable pour la coupure des mots en français
\usepackage[utf8]{inputenc} 
\usepackage{amsmath}
\usepackage{amssymb}
\usepackage{amsfonts}
\usepackage{stmaryrd}
\usepackage{xcolor}
\usepackage{titlesec}
\titleformat{\chapter}[block]
  {\normalfont\bfseries\filcenter}
  {}
  {0pt}
  {\titlerule\vspace{1ex}\Huge}
\usepackage[margin=1cm, includeheadfoot]{geometry}
\usepackage{frenchmath}

\begin{document}

\chapter{Mines-Ponts}


% --- TITRE DE L'EXERCICE (Style identique au titre 1.7 de l'image) ---
\begin{center}
    \large \textbf{1. Etude d'une matrice à diagonale nulle}
\end{center}
\vspace{0.2cm}

% --- ENCADRÉ DE L'ÉNONCÉ (Même style visuel avec source en bas à droite) ---
\noindent\framebox[\textwidth]{%
    \begin{minipage}{0.94\textwidth}
        \vspace{0.2cm}
        Soit un entier $n \geqslant 2$. On note $A$ la matrice de $\mathcal{M}_n(\mathbb{R})$ dont les coefficients diagonaux valent 0 et les autres coefficients valent 1.
        \vspace{0.2cm}
        
        \noindent $1.$ Calculer $A^2$. En déduire que $A$ est inversible et exprimer $A^{-1}$ ;
        
        \noindent $2.$ Déterminer les valeurs propres de $A$ et les espaces propres correspondants.
        \vspace{0.1cm}
        \begin{flushright}
            (\textbf{Mines-Ponts, PC (2025)})
        \end{flushright}
    \end{minipage}%
}
\vspace{0.4cm}


% --- TITRE DE L'EXERCICE ---
\begin{center}
    \large \textbf{2. Suites de Cauchy et complétude}
\end{center}
\vspace{0.2cm}

% --- ENCADRÉ DE L'ÉNONCÉ ---
\noindent\framebox[\textwidth]{%
    \begin{minipage}{0.94\textwidth}
        \vspace{0.2cm}
        Soit $(E, \| \cdot \|)$ un espace vectoriel normé. Une suite $(u_n) \in E^{\mathbb{N}}$ est de Cauchy si $\forall \varepsilon > 0, \exists N \in \mathbb{N}, \forall n, m \geqslant N, \|u_n - u_m\| \leqslant \varepsilon$.
        \vspace{0.2cm}
        
        \noindent \textbf{a)} Montrer que toute suite convergente est de Cauchy.
        
        \noindent \textbf{b)} Dans $E = \mathbb{R}[X]$ muni de la norme $\left\| \sum a_k X^k \right\| = \max |a_k|$, montrer que la suite $(P_n)$ de terme général $P_n = 1 + \sum_{k=1}^n \frac{X^k}{k}$ est de Cauchy sans être convergente.
        
        \noindent \textbf{c)} Montrer que toute suite de Cauchy est bornée.
        
        \noindent \textbf{d)} Montrer que, si $(u_n)$ est de Cauchy et possède une suite extraite convergente, alors $(u_n)$ est convergente.
        
        \noindent \textbf{e)} On admet le théorème de Bolzano-Weierstrass dans $\mathbb{R}$. Montrer que si $E$ est de dimension finie, alors la suite $(u_n)$ est convergente si et seulement si elle est de Cauchy.
        \vspace{0.1cm}
        \begin{flushright}
            (\textbf{Mines-Ponts, PSI (2025)})
        \end{flushright}
    \end{minipage}%
}
\vspace{0.4cm}
% --- TITRE DE L'EXERCICE ---
\begin{center}
    \large \textbf{3. Etude d'une relation d'anticommutation}
\end{center}

\vspace{0.2cm}
% --- ENCADRÉ DE L'ÉNONCÉ ---
\noindent\framebox[\textwidth]{%
\begin{minipage}{0.94\textwidth}
    \vspace{0.2cm}
    Soit $E$ un $\mathbb{K}$-espace vectoriel de dimension $n$. Soient $U$ et $V$ deux sous-espaces vectoriels de $\mathcal{L}(E)$ tels que $U + V = \mathcal{L}(E)$ et tels que $\forall u \in U, \forall v \in V, u \circ v + v \circ u = 0$.
    
    \vspace{0.2cm}
    \noindent \textbf{a)} Montrer qu'il existe $p \in U$ et $q \in V$ deux projecteurs tels que $p + q = \mathrm{id}$.
    
    \noindent \textbf{b)} Si $v \in V$, on pose $\varphi(v) = v_{|\operatorname{Ker}(p)} \in \mathcal{L}(\operatorname{Ker} p, E)$. Montrer que $\varphi$ est injective. En déduire que $\dim(V) \leqslant (n - r)^2$, avec $r = \operatorname{rg}(p)$.
    
    \noindent \textbf{c)} Montrer que $U$ ou $V$ est égal à $\{0\}$.
    \vspace{0.1cm}
    \begin{flushright}
            (\textbf{Mines-Ponts, PSI (2026) - RMS 860})
        \end{flushright}
\end{minipage}%
}
\vspace{0.4cm}
\pagebreak

% --- TITRE DE L'EXERCICE ---

\begin{center}

\large \textbf{4. Méthode probabiliste pour une inégalité de géometrie combinatoire}

\end{center}

\vspace{0.2cm}
% --- ENCADRÉ DE L'ÉNONCÉ ---

\noindent\framebox[\textwidth]{%

\begin{minipage}{0.94\textwidth}

\vspace{0.2cm}

Soient $E$ un espace euclidien et $(u_1, \dots, u_n)$ une famille de vecteurs unitaires de $E$. Soient $X_1, \dots, X_n$ des variables aléatoires indépendantes et de même loi uniforme sur $\{-1, 1\}$. Soit $X = \|X_1u_1 + \dots + X_nu_n\|^2$.

\vspace{0.2cm}
    \noindent \textbf{a)} Calculer $\mathbf{E}(X)$ ;
    
    \noindent \textbf{b)} Montrer qu'il existe $(\varepsilon_1, \dots, \varepsilon_n) \in \{-1, 1\}^n$ tel que $\|\varepsilon_1u_1 + \dots + \varepsilon_nu_n\| \leqslant \sqrt{n}$ .
    \vspace{0.1cm}
    \begin{flushright}
        \textbf{(Mines-Ponts, PSI (2026) - RMS 957)}
    \end{flushright}
\end{minipage}%
}

\vspace{0.4cm}
% --- TITRE DE L'EXERCICE ---
\begin{center}
    \large \textbf{5. Une équation fonctionnelle}
\end{center}
\vspace{0.2cm}

% --- ENCADRÉ DE L'ÉNONCÉ ---
\noindent\framebox[\textwidth]{%
    \begin{minipage}{0.94\textwidth}
        \vspace{0.2cm}
        Résoudre dans $\mathcal{C}^2(\mathbb{R}, \mathbb{R})$ l'équation fonctionnelle : $f(2x) = 4f(x) + 3x + 1$.
        \vspace{0.2cm}
        
        \begin{flushright}
            (\textbf{Mines-Ponts, PSI (2026) - RMS 910})
        \end{flushright}
    \end{minipage}%
}
\vspace{0.4cm}
\chapter{CentraleSupélec (Maths 1)}

% --- TITRE DE L'EXERCICE ---

\begin{center}

\large \textbf{1. Trace d'un endomorphisme des matrices}

\end{center}

\vspace{0.2cm}
% --- ENCADRÉ DE L'ÉNONCÉ ---

\noindent\framebox[\textwidth]{%

\begin{minipage}{0.94\textwidth}

\vspace{0.2cm}

    \noindent \textbf{a)} Montrer que : $\forall A, B \in \mathcal{M}_n(\mathbb{C}), \text{Tr}(AB) = \text{Tr}(BA)$.

    \noindent \textbf{b)} Soit $f \in \mathcal{L}(\mathcal{M}_n(\mathbb{C}), \mathcal{M}_p(\mathbb{C}))$ telle que : $\forall A, B \in \mathcal{M}_n(\mathbb{C}), f(AB) = f(A) f(B)$. Soit $\phi : M \in \mathcal{M}_n(\mathbb{C}) \mapsto \text{Tr}(f(M))$. Montrer : $\forall A, B \in \mathcal{M}_n(\mathbb{C}), \phi(AB) = \phi(BA)$ ;
    
    \noindent \textbf{c)} En déduire qu'il existe $\alpha \in \mathbb{C}$ tel que : $\forall M \in \mathcal{M}_n(\mathbb{C}), \text{Tr}(f(M)) = \alpha \text{Tr}(M)$ .
    \vspace{0.1cm}
    \begin{flushright}
        (\textbf{Centrale PSI (Maths 1) (2026) - RMS 1250})
    \end{flushright}
\end{minipage}%
}

\vspace{0.4cm}

\end{document}